\documentclass[11pt,a4paper]{article}

% =====================================================
% Packages
% =====================================================
\usepackage[utf8]{inputenc}
\usepackage[T1]{fontenc}
\usepackage[french]{babel}
\usepackage{geometry}
\usepackage{graphicx}
\usepackage{booktabs}
\usepackage{hyperref}
\usepackage{listings}
\usepackage{xcolor}
\usepackage{amsmath}
\usepackage{amssymb}
\usepackage{caption}
\usepackage{float}
\usepackage{enumitem}
\usepackage{array}

\geometry{margin=2.5cm}
\hypersetup{colorlinks=true, linkcolor=blue!50!black, urlcolor=blue!60!black}

% --- Style listings (code) ---
\definecolor{codegray}{rgb}{0.95,0.95,0.95}
\definecolor{codekw}{rgb}{0.13,0.29,0.53}
\definecolor{codestr}{rgb}{0.31,0.60,0.02}
\definecolor{codecom}{rgb}{0.4,0.4,0.4}
\lstset{
    backgroundcolor=\color{codegray},
    basicstyle=\ttfamily\footnotesize,
    breaklines=true,
    frame=single,
    numbers=left,
    numberstyle=\tiny\color{gray},
    keywordstyle=\color{codekw}\bfseries,
    stringstyle=\color{codestr},
    commentstyle=\color{codecom}\itshape,
    showstringspaces=false,
    language=Python,
    captionpos=b,
}

% =====================================================
% Page de titre
% =====================================================
\title{
    \vspace{-1.5cm}
    \textbf{Intelligence Artificielle}\\
    \large Projet : Variante déterministe du jeu de Domino 10\textbar{}10\\
    \vspace{0.3cm}
    \rule{0.6\textwidth}{0.5pt}
}
\author{
    \textsc{BENMEZIANE Samy} \\
    \texttt{samy.benmeziane@etu.u-paris.fr}
    \and
    \textsc{CHEGGOUR Racha} \\
    \texttt{racha.cheggour@etu.u-paris.fr}
}
\date{Mai 2026 \\ \small Université Paris Cité --- L3 Informatique}

\begin{document}
\maketitle

\begin{abstract}
Ce projet propose une variante déterministe et à information parfaite du jeu de Dominos
double-10 (66 dominos, 33 par joueur, distribution fixe), formalisée comme un problème
de recherche au sens du cours d'IA. Cinq intelligences artificielles ont été
implémentées : une IA aléatoire (baseline) et quatre IA à base de Minimax avec
élagage \(\alpha\)-\(\beta\), différenciées par leur fonction d'évaluation
(\texttt{f\_pip}, \texttt{f\_blocage}, \texttt{f\_protect}, \texttt{f\_mixte}) et leur
profondeur de recherche (4 ou 6). Un tournoi structuré de \textbf{600 parties}
(10 configurations stratifiées \(\times\) 3 variantes \(\times\) 2 attributions
\(\times\) 10 couples d'IA) a permis de mesurer les performances relatives,
l'avantage du premier joueur (\(\approx\) 61\%) et l'efficacité de l'élagage
\(\alpha\)-\(\beta\) (gain de \textbf{18.8\(\times\)} avec move ordering).
\end{abstract}

\tableofcontents
\newpage

% =====================================================
% Section 1 — Présentation du jeu (sujet 5.1)
% =====================================================
\section{Présentation générale du jeu}

\subsection{Choix du jeu et motivation}

Nous avons choisi une \textbf{variante du jeu de Dominos} dans laquelle nous éliminons
toute composante stochastique afin de respecter les contraintes du sujet :
\textit{jeu déterministe à information parfaite}. La variante classique du domino
implique une pioche cachée (talon) qui introduit du hasard et de l'information
incomplète~; nous l'avons supprimée.

Pour augmenter la complexité du jeu (facteur de branchement et profondeur de recherche
significatifs), nous travaillons avec un \textbf{jeu double-10} composé de
\textbf{66 dominos} (de \texttt{[0\textbar{}0]} à \texttt{[10\textbar{}10]}), distribués
\textbf{33 par joueur} dès l'initialisation. Cette taille a été choisie après que la
version classique double-6 (28 dominos) a été jugée trop simple par l'enseignante
encadrante.

\subsection{Règles retenues}

\begin{itemize}[leftmargin=*]
    \item Le joueur qui possède le \textbf{plus fort double} (10\textbar{}10 dans
          la quasi-totalité des cas) commence et \textbf{doit le poser} en premier coup.
    \item Le plateau est \textbf{linéaire} (extrémités gauche/droite, sans bifurcation
          aux doubles). Cette simplification limite la complexité tout en préservant
          la richesse stratégique.
    \item À chaque tour, un joueur pose un domino dont une des deux valeurs correspond
          à l'une des deux extrémités du plateau. S'il ne le peut pas, il \textbf{passe
          son tour} (sans en être pénalisé directement).
    \item Le jeu s'arrête lorsque (a) un joueur a vidé sa main, ou (b) les deux joueurs
          sont simultanément bloqués.
    \item En cas de blocage, le \textbf{joueur ayant le moins de pips} restants gagne.
\end{itemize}

\subsection{Modélisation comme problème de recherche}

Nous formalisons le jeu en suivant les définitions du cours :

\begin{itemize}[leftmargin=*]
    \item \textbf{État initial $s_0$} : une distribution fixe et reproductible des
          66 dominos en deux mains (33 chacune), un plateau vide, et un joueur
          courant déterminé par le possesseur du plus fort double.
    \item \textbf{Actions(s)} : ensemble des couples (domino, position) légaux pour
          le joueur courant.
    \item \textbf{Result(s, a)} : nouvel état après pose du domino, avec mise à jour
          des extrémités du plateau, retrait du domino de la main, et changement de
          joueur courant.
    \item \textbf{TerminalTest(s)} : $s$ terminal si une main est vide ou si les
          deux joueurs sont bloqués.
    \item \textbf{Utility(s)} : $+1$ pour MAX en victoire, $-1$ en défaite, $0$ en
          égalité (jeu à somme nulle).
\end{itemize}

% =====================================================
% Section 2 — Moteur de jeu (sujet 5.2)
% =====================================================
\section{Implémentation et validation du moteur de jeu}

\subsection{Architecture orientée objet}

L'implémentation est en \textbf{Python 3.10+} (sans dépendance plateforme : compatible
Linux et Windows). L'architecture suit une décomposition orientée objet stricte :

\begin{table}[H]
\centering
\begin{tabular}{ll}
\toprule
\textbf{Module} & \textbf{Rôle} \\
\midrule
\texttt{src/domino.py} & Classe \texttt{Domino} immuable (\texttt{\_\_slots\_\_}), génération du jeu complet \\
\texttt{src/plateau.py} & Plateau linéaire avec extrémités, génération des coups légaux \\
\texttt{src/joueur.py} & Joueur abstrait + \texttt{JoueurHumain} (CLI) \\
\texttt{src/partie.py} & Moteur de jeu : état, transitions, terminaison, utility \\
\texttt{src/ia/} & Module des IA : aléatoire, Minimax, fonctions d'évaluation \\
\texttt{src/tournoi/} & Génération de configurations et runner de tournoi \\
\texttt{tests/} & 35 tests unitaires (pytest) \\
\bottomrule
\end{tabular}
\caption{Organisation des modules du projet.}
\end{table}

\subsection{Choix d'implémentation clés}

\paragraph{Domino immuable.} La classe \texttt{Domino} utilise \texttt{\_\_slots\_\_}
et stocke les valeurs $(a, b)$ avec la convention $a \leq b$. Elle est \textbf{strictement
immuable}, ce qui permet de partager les références sans risque de mutation accidentelle.

\paragraph{Clone léger.} Un point critique de performance pour le Minimax était le
clonage de la \texttt{Partie} pour explorer les coups. Notre première implémentation
utilisait \texttt{copy.deepcopy()} qui se révélait extrêmement lent (de l'ordre de
60$\times$ plus lent que nécessaire). Nous avons réécrit un \textbf{clone léger}
exploitant l'immuabilité des dominos : seules les listes mutables sont copiées
explicitement, les références aux dominos étant partagées sans risque.

\paragraph{Premier coup forcé.} La règle exigeant que le possesseur du plus fort
double le pose en premier est gérée au niveau de la classe \texttt{Partie} via la
méthode \texttt{coup\_force()}, transparente pour les IA.

\subsection{Validation par tests unitaires}

Le projet contient \textbf{35 tests unitaires} (avec \texttt{pytest}), tous verts en
moins de 10 secondes. Ils couvrent :

\begin{itemize}[leftmargin=*]
    \item \textbf{Domino} (7 tests) : création normalisée, calcul des pips, doubles, contient/autre côté, égalité/hash, validation des bornes, génération du jeu complet (66 dominos).
    \item \textbf{Plateau} (7 tests) : état initial vide, pose initiale, coups possibles, pose à gauche/droite, refus des coups illégaux, indépendance des simulations.
    \item \textbf{Partie} (8 tests) : détermination du premier joueur, coup forcé, victoire, blocage, et 3 tests d'\textbf{indépendance du clone léger} (mains, plateau, joueur courant).
    \item \textbf{IA et fonctions d'évaluation} (13 tests) : chaque niveau d'IA produit un coup légal, parties complètes, comportement des 4 fonctions d'évaluation incluant les 3 cas hiérarchisés de \texttt{f\_blocage}.
\end{itemize}

\includegraphics[width=0.95\textwidth]{todo1.png}

% =====================================================
% Section 3 — Évaluations + alpha-beta (sujet 5.3)
% =====================================================
\section{Fonctions d'évaluation, score et élagage \(\alpha\)-\(\beta\)}

\subsection{Algorithme Minimax avec élagage \(\alpha\)-\(\beta\)}

Toutes nos IA "intelligentes" reposent sur l'algorithme \textbf{Minimax} avec
\textbf{élagage \(\alpha\)-\(\beta\)} et \textbf{move ordering} (tri des coups par
pips décroissants pour favoriser les coupures précoces). L'évaluation aux feuilles
utilise une des quatre fonctions décrites ci-dessous, et un état terminal renvoie
$\pm 10\,000$ pour écraser toute heuristique et garantir la priorité absolue à la
victoire.

\subsection{Description des fonctions d'évaluation}

Nous avons développé \textbf{quatre fonctions d'évaluation distinctes}, chacune
incarnant une stratégie différente. Toutes sont calculées du point de vue du joueur
MAX. Soit $g$ et $d$ les valeurs aux extrémités gauche et droite du plateau~;
$\text{pips}(j)$ la somme des pips des dominos en main du joueur $j$.

\paragraph{(1) \texttt{f\_pip} --- différentiel de pips.}
\[
    f_{\text{pip}}(s) = \text{pips}(\text{MIN}) - \text{pips}(\text{MAX})
\]
\textit{Intuition :} récompense les positions où l'adversaire garde beaucoup de pips
et où nous en avons disposé. C'est la stratégie la plus directe : se rapprocher de
l'épuisement de sa propre main.\\
\textit{Bornes :} $\left[-\text{pips}_0(\text{MAX}),\, +\text{pips}_0(\text{MIN})\right]$
où $\text{pips}_0$ représente la somme initiale.

\paragraph{(2) \texttt{f\_blocage} --- pénaliser l'adversaire bloqué (3 cas hiérarchisés).}
Cette fonction implémente une \textbf{hiérarchie explicite de cas} :
\begin{itemize}[leftmargin=*]
    \item \textbf{Cas 1 (blocage total)} : aucun domino adverse ne contient ni $g$ ni $d$.
          \[ f_{\text{blocage}} = \text{pips}(\text{MIN}) \]
    \item \textbf{Cas 2 (blocage partiel)} : une seule extrémité ne contient aucune valeur en main adverse.
          \[ f_{\text{blocage}} = \text{pips}(\text{MIN}) - 1 \]
    \item \textbf{Cas 3 (aucune extrémité bloquée)} :
          \[ f_{\text{blocage}} = \sum_{d \in \text{main}_{\text{MIN}},\ d \not\ni g \wedge d \not\ni d}\!\!\!\!\!\! \text{pips}(d) \]
\end{itemize}
\textit{Intuition :} privilégier strictement le blocage total sur le partiel
(différence garantie de 1 point), tout en récompensant graduellement les blocages
individuels de dominos.\\
\textit{Bornes :} $\left[0, \text{pips}_0(\text{MIN})\right]$.

\paragraph{(3) \texttt{f\_protect} --- éviter d'être soi-même bloqué (miroir négatif).}
Strict miroir de \texttt{f\_blocage} appliqué à notre propre main, avec signe négatif :
\begin{itemize}[leftmargin=*]
    \item Cas 1 (totalement bloqué) : $f_{\text{protect}} = -\text{pips}(\text{MAX})$
    \item Cas 2 (partiellement) : $f_{\text{protect}} = -(\text{pips}(\text{MAX}) - 1)$
    \item Cas 3 (libre) : $f_{\text{protect}} = -\!\!\sum \text{pips}$ des dominos isolés en main MAX
\end{itemize}
\textit{Bornes :} $\left[-\text{pips}_0(\text{MAX}), 0\right]$.

\paragraph{(4) \texttt{f\_mixte} --- somme des trois.}
\[
    f_{\text{mixte}} = f_{\text{pip}} + f_{\text{blocage}} + f_{\text{protect}}
\]
Combine course aux pips, attaque (blocage de l'adversaire) et défense (anti-blocage de soi).

\subsection{Analyse expérimentale de l'élagage \(\alpha\)-\(\beta\)}

Nous avons mesuré l'impact concret de l'élagage \(\alpha\)-\(\beta\) sur une partie
complète à profondeur 4 avec \texttt{f\_pip}.

\begin{table}[H]
\centering
\begin{tabular}{lrrr}
\toprule
\textbf{Configuration} & \textbf{Temps} & \textbf{Nœuds explorés} & \textbf{Élagages} \\
\midrule
Minimax brut (sans optim.) & 3.39 s & 171\,732 & 0 \\
Minimax + \(\alpha\)-\(\beta\) & 0.67 s & 31\,064 & 5\,394 \\
Minimax + \(\alpha\)-\(\beta\) + move ordering & \textbf{0.18 s} & \textbf{9\,634} & 2\,781 \\
\bottomrule
\end{tabular}
\caption{Impact mesuré de l'élagage \(\alpha\)-\(\beta\) et du move ordering
(profondeur 4, \texttt{f\_pip}, partie complète).}
\label{tab:alphabeta}
\end{table}

\begin{figure}[H]
\centering
\includegraphics[width=0.95\textwidth]{impact_alpha_beta.png}
\caption{Impact de \(\alpha\)-\(\beta\) et move ordering sur le temps, le nombre de
nœuds explorés et le nombre d'élagages déclenchés.}
\label{fig:alphabeta}
\end{figure}

\textbf{Lectures principales :}
\begin{itemize}[leftmargin=*]
    \item L'élagage \(\alpha\)-\(\beta\) seul divise par \textbf{5} le temps et par
          \textbf{5.5} le nombre de nœuds explorés, conformément à la théorie.
    \item Le \textbf{move ordering} apporte un facteur \(3.7\)$\times$ supplémentaire,
          aboutissant à un gain total de \textbf{18.8$\times$}.
    \item Paradoxalement, le nombre d'élagages \emph{diminue} avec move ordering :
          en explorant les coups les plus prometteurs en premier, on déclenche les
          coupures plus tôt, ce qui réduit l'arbre exploré et donc le nombre
          d'opportunités de coupures.
    \item Le résultat de la partie est strictement \textbf{identique} dans les trois
          configurations : \(\alpha\)-\(\beta\) ne modifie pas la qualité du jeu,
          uniquement la vitesse de calcul.
\end{itemize}

  \begin{figure}[H]
  \centering
  \includegraphics[width=0.95\textwidth]{noeuds_vs_profondeur.png}
  \caption{Croissance du nombre de nœuds explorés et du temps de calcul en
  fonction de la profondeur Minimax (échelle logarithmique). Trois configurations
  sont comparées~: Minimax seul (rouge), avec \(\alpha\)-\(\beta\) (orange), et avec
  \(\alpha\)-\(\beta\) + move ordering (vert). Les optimisations s'\textbf{empilent
  de manière multiplicative}~: à profondeur 6, le gain total atteint
  \textbf{200\(\times\)}.}
  \label{fig:profondeur}
  \end{figure}

  \begin{table}[H]
  \centering
  \small
  \begin{tabular}{rrrrrrr}
  \toprule
  \textbf{Prof.} & \multicolumn{2}{c}{\textbf{Sans \(\alpha\)-\(\beta\)}} &
  \multicolumn{2}{c}{\textbf{Avec \(\alpha\)-\(\beta\) seul}} &
  \multicolumn{2}{c}{\textbf{Avec \(\alpha\)-\(\beta\) + ordering}} \\
  \cmidrule(lr){2-3}\cmidrule(lr){4-5}\cmidrule(lr){6-7}
         & Nœuds       & Temps   & Nœuds    & Temps   & Nœuds  & Temps   \\
  \midrule
  2  &       98 & 0.002 s &      66 & 0.001 s &      28 & 0.001 s \\
  3  &      950 & 0.021 s &     353 & 0.008 s &     134 & 0.003 s \\
  4  &    9\,626 & 0.208 s &  1\,193 & 0.026 s &     544 & 0.012 s \\
  5  &   93\,996 & 2.079 s &  6\,972 & 0.148 s &  1\,959 & 0.044 s \\
  6  &  904\,830 & 19.6 s  & 34\,099 & 0.72 s  &  4\,519 & 0.10 s  \\
  \bottomrule
  \end{tabular}
  \caption{Mesures détaillées de l'impact de \(\alpha\)-\(\beta\) et du move
  ordering en fonction de la profondeur (1 coup d'IA après le 10\textbar{}10 forcé,
  config \texttt{ouverte\_01}/\texttt{v0}, eval \texttt{f\_pip}).}
  \label{tab:profondeur}
  \end{table}

% =====================================================
% Section 4 — Tournoi (sujet 5.4)
% =====================================================
\section{Tournoi entre les intelligences artificielles}

\subsection{Méthodologie : configurations stratifiées}

Le caractère strictement déterministe du jeu nous a posé un défi méthodologique :
deux IA déterministes affrontant la même configuration produisent toujours la même
partie. Pour respecter l'exigence de "$\geq 50$ parties par couple" tout en
disposant de données statistiquement variées, nous avons mis en place une
\textbf{méthodologie de tournoi stratifié} :

\begin{itemize}[leftmargin=*]
    \item \textbf{5 types de configurations} initiales définies par des critères
          quantitatifs sur la diversité de valeurs et le différentiel de pips :
          \texttt{ouverte}, \texttt{fermée}, \texttt{équilibrée},
          \texttt{avantage\_alpha}, \texttt{avantage\_beta}.
    \item \textbf{2 configurations par type}, soit 10 configurations totales,
          générées de manière reproductible par un seed fixe (42).
    \item \textbf{3 variantes par configuration} (notées \texttt{v0}, \texttt{v1},
          \texttt{v2}) : la \texttt{v0} est l'originale, et \texttt{v1}/\texttt{v2}
          en dérivent par échange aléatoire de 2 dominos en préservant le type
          de configuration.
    \item \textbf{2 attributions des mains} par variante : la même configuration est
          jouée une fois avec l'IA-A possédant \texttt{Main\_alpha} (contenant le
          10\textbar{}10 et donc starter), et une fois avec l'IA-B la possédant.
\end{itemize}

\textbf{Total :} $10 \times 3 \times 2 = 60$ parties par couple d'IA, soit \textbf{600 parties}
en additionnant les 10 couples uniques formés à partir des 5 IA.

\subsection{Cinq IA, quatre fonctions d'évaluation}

\begin{table}[H]
\centering
\begin{tabular}{lcccl}
\toprule
\textbf{IA} & \textbf{Profondeur} & \textbf{$\alpha$-$\beta$} & \textbf{Move ordering} & \textbf{Évaluation} \\
\midrule
\texttt{IA1\_aleatoire} & --- & --- & --- & hasard \\
\texttt{IA2\_pip} & 6 & \checkmark & \checkmark & \texttt{f\_pip} \\
\texttt{IA3\_blocage} & 6 & \checkmark & \checkmark & \texttt{f\_blocage} \\
\texttt{IA4\_protect} & 6 & \checkmark & \checkmark & \texttt{f\_protect} \\
\texttt{IA5\_mixte} & 4 & \checkmark & \checkmark & \texttt{f\_mixte} \\
\bottomrule
\end{tabular}
\caption{Configuration des 5 IA en lice. La profondeur réduite de \texttt{IA5\_mixte}
(4 contre 6 pour les autres) permet de tester si une stratégie multi-critères mais peu
profonde peut compenser un manque de prévision face à des stratégies mono-critères
mais plus profondes.}
\end{table}

\subsection{Résultats globaux}

\begin{figure}[H]
\centering
\includegraphics[width=0.85\textwidth]{classement_global.png}
\caption{Classement global des 5 IA par taux de victoire (toutes oppositions confondues),
avec intervalles de confiance à 95\%. L'aléatoire (IA1) est nettement isolée à 13.3\%~;
les quatre IA Minimax sont regroupées entre 55\% et 60\% avec des intervalles de
confiance qui se chevauchent largement.}
\label{fig:classement}
\end{figure}

\begin{table}[H]
\centering
\begin{tabular}{lcc}
\toprule
\textbf{IA} & \textbf{Victoires / 240} & \textbf{Taux} \\
\midrule
\texttt{IA2\_pip}       & 145 & \textbf{60.4\%} \\
\texttt{IA3\_blocage}   & 143 & 59.6\% \\
\texttt{IA4\_protect}   & 139 & 57.9\% \\
\texttt{IA5\_mixte}     & 133 & 55.4\% \\
\texttt{IA1\_aleatoire} &  32 & 13.3\% \\
\bottomrule
\end{tabular}
\caption{Classement global après 600 parties au total (240 par IA).}
\end{table}

\subsection{Matrice détaillée des affrontements}

\begin{figure}[H]
\centering
\includegraphics[width=0.85\textwidth]{significativite_statistique.png}
\caption{Matrice de domination 5\(\times\)5 (taux de victoire de A contre B). Les
\textbf{cases hachurées} indiquent les couples dont l'écart de victoire n'est
\emph{pas} statistiquement significatif au seuil de 95\% (intervalle de confiance
$\pm 13\%$ à $n=60$).}
\label{fig:heatmap_signif}
\end{figure}

\textbf{Observations :}
\begin{itemize}[leftmargin=*]
    \item Les couples impliquant \texttt{IA1\_aleatoire} sont tous \emph{très
          significatifs} : toutes les IA Minimax la battent à 85--88\%, confirmant
          l'apport décisif de la planification.
    \item En revanche, parmi les 6 couples Minimax-vs-Minimax, \textbf{seuls 1 ou 2
          sont statistiquement significatifs} :
          \begin{itemize}
              \item \texttt{IA3\_blocage} > \texttt{IA4\_protect} (58.3\% / 40.0\%) : significatif.
              \item \texttt{IA2\_pip} > \texttt{IA5\_mixte} (56.7\% / 41.7\%) : significatif.
              \item Les autres écarts (\texttt{IA2\_pip} vs \texttt{IA3\_blocage}, etc.) tombent dans le bruit statistique.
          \end{itemize}
\end{itemize}

\subsection{Triangle non-transitif (rock-paper-scissors)}

\begin{figure}[H]
\centering
\includegraphics[width=0.65\textwidth]{triangle_non_transitif.png}
\caption{Relations de domination entre les 3 stratégies pures. Une seule des trois
flèches du triangle est statistiquement significative (\texttt{blocage} > \texttt{protect}).}
\label{fig:triangle}
\end{figure}

Les données suggèrent un triangle "rock-paper-scissors" :
\(\texttt{pip} \succ \texttt{blocage} \succ \texttt{protect} \succ \texttt{pip}\).
\textbf{Néanmoins, deux des trois relations ne sont pas statistiquement significatives}
(écarts de seulement 1.7~point à 60 parties). Le triangle est donc \emph{partiellement
illusoire} et son interprétation doit être nuancée. Pour le confirmer, un tournoi
plus volumineux ($n \geq 200$ parties par couple, IC $\pm 7\%$) serait nécessaire.

\subsection{Avantage du premier joueur}

\begin{figure}[H]
\centering
\includegraphics[width=0.85\textwidth]{avantage_premier_joueur.png}
\caption{Taux de victoire du joueur qui commence (le \emph{starter}), couple par
couple. Tous les couples sont au-dessus de la barre des 50\% (rouge), avec une
moyenne globale de 61.2\% (verte).}
\label{fig:starter}
\end{figure}

\textbf{Avantage starter mesuré : 61.2\% en moyenne}, avec une plage entre 53.3\% et
66.7\% selon les couples. Cet avantage est mécaniquement explicable : le starter pose
le 10\textbar{}10 dès son premier coup, soit une réduction immédiate de 20 pips
(c.~à~d. environ 6\% de la masse totale de pips initiale) que l'adversaire ne peut
pas obtenir d'un coup unique.

\subsection{Performance par type de configuration}

\begin{figure}[H]
\centering
\includegraphics[width=0.95\textwidth]{performance_par_type_config.png}
\caption{Taux de victoire par IA en fonction du type de configuration initiale.}
\label{fig:par_type}
\end{figure}

\textbf{Résultat marquant} : sur les configurations de type \texttt{fermee} (où une
main est volontairement concentrée sur peu de valeurs, créant un risque de blocage),
\textbf{toutes les IA atteignent un taux de victoire de 50\%} -- y compris l'IA
aléatoire. Cela s'explique par la \emph{symétrie parfaite} de notre design~: pour
chaque variante, nous jouons la même partie deux fois en inversant l'attribution des
mains, ce qui annule statistiquement l'effet de la stratégie sur ce type de
configuration. Cela suggère que dans les positions très fermées, le jeu est
\textbf{déterminé par la distribution initiale}, et la stratégie n'a plus d'influence.

\subsection{Coût temporel}

\begin{figure}[H]
\centering
\includegraphics[width=0.95\textwidth]{duree_moyenne.png}
\caption{Temps moyen par partie pour chaque couple d'IA (échelle logarithmique).}
\label{fig:duree}
\end{figure}

Le temps total du tournoi a été de \textbf{27 min 35 s} (1655 s) pour 600 parties.
Sans \(\alpha\)-\(\beta\), une extrapolation indique qu'il aurait fallu environ
\textbf{8 heures}.

% =====================================================
% Section 5 — IA générative (sujet 5.5)
% =====================================================
 \section{Utilisation de l'IA générative}

Conformément à la consigne, nous précisons ici notre usage ponctuel d'outils d'IA
générative durant ce projet. La conception, les choix algorithmiques, la
méthodologie expérimentale et l'analyse des résultats sont entièrement notre
travail ; l'IA a uniquement été sollicitée comme assistance technique.

\begin{description}[leftmargin=*]
    \item[Modèle utilisé.] Assistance ponctuelle via une IA générative
    conversationnelle.

    \item[Type d'utilisation.] Aide ponctuelle pour la clarification de certains
    points de syntaxe Python ou LaTeX, la reformulation de certaines phrases du
    rapport, la vérification de petits points techniques et le débogage occasionnel
    de messages d'erreur.

    \item[Environnement.] Navigateur Web, en parallèle de l'éditeur principal
    VS Code sous Windows.

    \item[Rôle effectif.] L'IA a joué un rôle d'assistance ciblée. Nous lui avons
    posé des questions précises, puis nous avons vérifié et adapté les réponses
    avant toute intégration. La conception des fonctions d'évaluation, la
    méthodologie du tournoi, l'implémentation du clone léger pour Minimax et
    l'analyse statistique des résultats restent notre travail.

    \item[Apport et limites.] L'IA a permis de gagner du temps sur quelques
    ambiguïtés syntaxiques et rédactionnelles. Cependant, certaines réponses
    étaient incomplètes ou incohérentes avec notre architecture, ce qui a nécessité
    un regard critique systématique.
\end{description}
% =====================================================
% Section 6 — Difficultés rencontrées (sujet 5.6)
% =====================================================
\section{Difficultés rencontrées}\label{sec:difficultes}

\paragraph{(1) Performance du Minimax.} La première implémentation utilisait
\texttt{copy.deepcopy()} pour cloner l'état lors de l'exploration de l'arbre de
recherche. Mesurée à \textbf{60$\times$ trop lente}, cette approche rendait un
tournoi à profondeur 6 totalement inenvisageable (estimation de plus de 8 heures).
Nous avons réécrit un \emph{clone léger} exploitant l'immuabilité des dominos~:
seules les listes mutables (mains, plateau) sont copiées explicitement. Ce
changement a réduit le temps du tournoi à \textbf{27 minutes}.

\paragraph{(2) Non-déterminisme du tournoi.} Notre jeu étant déterministe, jouer
50 fois la même configuration produisait 50 fois le même résultat. Nous avons
résolu ce paradoxe en adoptant une \emph{méthodologie de stratification}
(5 types $\times$ 2 configurations $\times$ 3 variantes $\times$ 2 attributions) qui
garantit la diversité des situations testées sans introduire d'aléatoire.

\paragraph{(3) Calibration des fonctions d'évaluation.} Plusieurs versions
intermédiaires de \texttt{f\_blocage} ont été testées avant d'arriver à la version
finale à 3 cas hiérarchisés. Une version trop agressive sur le blocage
(\(+200\) pour blocage total, \(+50\) pour quasi-blocage) a même \emph{dégradé}
les performances par rapport à une version simple : trop de bonus rares et
non-linéaires faisait que l'IA prenait des décisions pip-suboptimales dans la
majorité des positions où aucun blocage n'était imminent. La version finale
(récompense proportionnelle à la masse de pips bloquée chez l'adversaire) est plus
contextuelle et performe mieux.

\paragraph{(4) Bug d'interface graphique.} Lors du refactoring de la GUI Tkinter
pour ajouter une scrollbar à la main du joueur, un \texttt{return} mal placé a
laissé tout le code des boutons et de l'historique \emph{après} le retour de la
fonction utilitaire \texttt{\_construire\_main\_scrollable}, donc \emph{jamais
exécuté}. Le bug s'est manifesté par des fenêtres sans aucun bouton et un plateau
qui restait vide. Détecté visuellement par l'utilisateur lors d'un test, il a été
corrigé en déplaçant le code de construction des contrôles dans la méthode
appropriée.

\paragraph{(5) Significativité statistique.} Avec 60 parties par couple,
l'intervalle de confiance à 95\% pour un taux de victoire est d'environ
$\pm 13\%$. La majorité des écarts entre IA Minimax tombent dans cette plage,
ce qui nous a contraints à présenter les résultats avec nuance plutôt que de
proclamer des vainqueurs nets. C'est un rappel utile que toute conclusion
doit être assortie d'une mesure de confiance.

% =====================================================
% Conclusion (optionnelle, courte)
% =====================================================
\section*{Conclusion}



Ce projet nous a permis d’implémenter et de comparer rigoureusement plusieurs stratégies
d’IA pour une variante déterministe du jeu de Dominos. Au-delà des résultats numériques (l’IA
\texttt{pip} légèrement en tête à 60.4%, l’avantage du starter mesuré à 61.2%, l’élagage
\(\alpha\)-\(\beta\) accélérant le calcul d’un facteur 18.8$\times$), l’apport principal est
méthodologique : \textbf{stratification du tournoi}, \textbf{tests d’indépendance du clone
Minimax}, et \textbf{analyse des intervalles de confiance}.

Un point notable est la performance légèrement supérieure de l’IA \texttt{pip}, qui s’est révélée
surprenante, dans la mesure où nous nous attendions initialement à ce que l’IA \texttt{Mix} soit
la plus performante. Une explication possible réside dans la différence de profondeur de recherche :
l’IA \texttt{Mix} est limitée à une profondeur de 4, tandis que les autres stratégies sont évaluées
avec une profondeur de 6. Malgré cela, ses performances restent proches, ce qui suggère que la
combinaison d’heuristiques peut partiellement compenser une profondeur plus faible.

Cependant, plusieurs résultats restent à confirmer sur un tournoi plus volumineux (notamment
l’existence ou non d’un triangle non-transitif). Le nombre de parties réalisées demeure relativement
limité, ce qui peut introduire une variabilité statistique. Une extension naturelle de ce travail
consisterait à augmenter significativement le nombre de simulations (par exemple, plusieurs centaines
de parties par confrontation), afin d’obtenir des conclusions plus robustes et généralisables.


\end{document}

